\section{Memento}
\label{sec:pat-memento}

\problemstatement{Capture and externalise an object's internal state
without violating encapsulation, so the object can be restored to that
state later -- the mechanism behind save points and undo.}

\begin{patternbox}[When You Reach for It]
The trigger is \emph{``save and restore an object's state''} -- snapshots,
checkpoints, undo history. When you must be able to roll an object back but
do not want to expose its internals to the code that stores the snapshot,
use a Memento.
\end{patternbox}

\subsection*{Understanding the pattern}

Three roles: the \textbf{Originator} (whose state we snapshot) creates a
\textbf{Memento} holding a copy of its state and can later restore from
one; the \textbf{Caretaker} keeps mementos (e.g.\ on an undo stack) but
treats them as opaque -- it never reads their contents. Encapsulation is
preserved because only the Originator can read or write a memento's
internals.

\begin{ideabox}[Opaque Snapshots the Caretaker Cannot Read]
The Originator packages its own state into a memento and unpacks it on
restore; the Caretaker only stores and hands mementos back. This division
gives you undo without leaking private state to the code that manages
history.
\end{ideabox}

\subsection*{C++ implementation}

\begin{lstlisting}
#include <bits/stdc++.h>
using namespace std;

class Memento {                                 // opaque snapshot
    friend class Editor;                        // only Editor sees inside
    string state;
    Memento(string s) : state(move(s)) {}
};

class Editor {                                  // originator
    string content;
public:
    void type(const string& text) { content += text; }
    string show() const { return content; }
    Memento save() const { return Memento(content); }      // snapshot
    void restore(const Memento& m) { content = m.state; }  // roll back
};

int main() {
    Editor editor;
    editor.type("Hello");
    Memento checkpoint = editor.save();         // caretaker holds this
    editor.type(", World");
    cout << "Before undo: " << editor.show() << "\n";
    editor.restore(checkpoint);                 // back to "Hello"
    cout << "After undo:  " << editor.show() << "\n";
    return 0;
}
\end{lstlisting}

\begin{complexitybox}[Trade-offs]
\textbf{Pros.} Clean save/restore without breaking encapsulation; pairs
naturally with Command for undo. \textbf{Cons.} Snapshots can be
memory-heavy if state is large or saved often; deciding what to include in
a memento and how frequently to snapshot needs care (consider
incremental/diff mementos).
\end{complexitybox}

\begin{takeawaybox}
Memento snapshots an object's state into an opaque object that only the
originator can interpret, enabling undo and checkpoints while keeping
internals private. It frequently teams with Command: the command stores a
memento to reverse itself. Cost to flag: memory for many/large snapshots.
\end{takeawaybox}
